% Conceptual: let-binding is NOT the same as textual substitution under CBV.
% Two independent failure modes, both verified in SML/NJ:
%   - x unused, e1 raises:  let raises, substitution gives 5
%   - x used twice, e1 prints: let prints P once, substitution prints PP
%     (and both still return 2 -- the values agree, only the effects differ,
%      which is why "same answer" is not the same as "equivalent")
% Safe exactly when e1 is a VALUE. Also verified: x used exactly once is NOT
% sufficient (order of effects can still shift in general), so the clean
% condition to grade on is "e1 is a value" / effect-free.
\part[3]\hfill

  True or false: \code{let val x = e1 in e2 end} is equivalent to
  $[\code{e1}/x]\,\code{e2}$, the substitution of \code{e1} for all occurrences
  of \code{x} in \code{e2}. If false, give a counterexample.

  \begin{solutionorbox}[15em]
    \textbf{False.} \code{let} evaluates \code{e1} \textbf{once, up front};
    substitution instead copies the unevaluated expression to each occurrence of
    \code{x} --- so the number of times \code{e1} runs can change.

    Take \code{e1} $=$ \code{(print "P"; 1)} and \code{e2} $=$ \code{x + x}. The
    \code{let} prints \code{P} once; the substituted form
    \code{(print "P"; 1) + (print "P"; 1)} prints it \textbf{twice}. Both return
    \code{2}, so this is a case where the two agree on the value and still are
    not equivalent.

    Taking \code{e2} $=$ \code{5}, with \code{x} unused, breaks it the other
    way: for \code{e1} $=$ \code{raise Fail ""} the \code{let} raises, while the
    substituted form is just \code{5}, which evaluates fine. Here \code{let}
    runs \code{e1} once and substitution runs it \textbf{zero} times.

    The claim holds when \code{e1} is a \textbf{value} --- there is then nothing
    to evaluate, so duplicating or dropping it costs nothing. (More generally it
    suffices that \code{e1} be effect-free and terminating.)
  \end{solutionorbox}
